\chapter{绪论}
\label{chap:introduction}
\section{研究背景与意义}
本章给出绪论的组织示例。撰写时应从实际需求或学科发展出发，逐步说明研究对象、关键问题及研究意义。模板中的文字用于指导结构安排，不代表任何真实研究结论，也不应直接作为正式论文内容提交。

建议先说明问题发生的条件，再解释现有方法难以满足哪些需求，最后明确本文研究的必要性。背景介绍应服务于研究问题，避免堆积与主题关联较弱的材料。引用他人成果时应在相应位置给出文献来源，例如文献\cite{example-book,example-article}。本模板的参考文献均明确标为虚构示例，正式使用时必须替换。
\section{国内外研究现状}
\subsection{研究方向一}
可按照方法类别、问题层次或发展阶段组织研究现状。每组文献不仅要介绍已有工作，还应比较其基本假设、适用条件和局限，避免把综述写成互不关联的文献清单。
\subsection{研究方向二}
围绕与论文核心问题直接相关的研究展开讨论。区分文献已报告的事实与作者自己的分析；对于有争议的判断，应说明比较依据和适用范围。
\subsubsection{研究现状的归纳方式}
这里展示第四级标题。通常不必为了细分而增加标题层级，能够用连续段落说清楚的内容可直接写在上一级标题下。
\section{研究内容与论文结构}
本文的占位研究内容包括问题建模、方法设计和结果分析。实际写作时，应说明各项工作之间的联系，而非简单重复各章标题。

第\ref{chap:examples}章展示公式、插图、表格和算法的写法；第\ref{chap:results}章展示结果分析的组织方式；最后一章给出总结与展望的写作提示。章节编号由交叉引用自动生成，调整章节顺序后无需手动改号。
\section{本章小结}
本章小结应归纳本章得到的认识，并说明它如何引出后续研究。避免将小结写成标题的机械复述。
